\documentclass[../../../../main.tex]{subfiles}
\begin{document}

\begin{flushright}
\footnotesize\textit{【自動文字起こし・要確認】}
\end{flushright}

[解]
(1) $N=5$の時
$i=j=0$での成立が必要だが
\begin{align*}
&P(X=0 \land Y=0) = \frac{1}{36} \\
&P(X=0) = \frac{1}{6} \quad (N=6\text{の時の図}) \\
&P(Y=0) = \frac{7}{36}
\end{align*}
から $P(X=0 \land Y=0) = P(X=0)P(Y=0)$ は不成立、よって$X, Y$は独立ではない。

$2^{\circ} \ N=6$の時
確率変数$X, Y$に対し、まず$P(X=k)=P(Y=k)=\frac{1}{6} \quad (k=0 \sim 5)$であり
$$ P(X=i \land Y=j) = \frac{1}{36} \quad \left( \text{$X=i$になるし(確率$1/6$)、$Y=j$となる目が出る(確率$1/6$)} \right) $$
から、任意の$i, j$に対し $P(X=i \land Y=j) = P(X=i)P(Y=j)$ が成立するので、$X$と$Y$は独立。

(2) $i=j=0$での成立が必要である。
$1^{\circ} \ N=3$の時
$P(X=k)=P(Y=k)=\frac{1}{3} \quad (k=0 \sim 2)$ であり (1)と同じく、
$P(X=i \land Y=j) = \frac{1}{9}$ だから、$X, Y$は独立。

$2^{\circ} \ N=4$の時、
$P(X=0) = \frac{1}{6}, \quad P(Y=0) = \frac{9}{36}, \quad P(X=0 \land Y=0) = \frac{1}{36}$ から
$X, Y$は非独立。

$3^{\circ} \ 7 \le N$の時
2週まわって6にいくことはないので、1回目に出た目を$A$、2回目を$B$とすると、$Y=6$となるのは $A+B=6$ の5通り。
$$ P(Y=6) = \frac{5}{36} $$
又、$P(X=6)=\frac{1}{6}, \quad P(X=6 \land Y=6)=0$ だから
$$ P(X=6)P(Y=6) \neq P(X=6 \land Y=6) $$
となり、$X$と$Y$は独立ではない。

以上から、$X, Y$が独立なのは
$$ N=3, 6 $$
の時。

\end{document}
